%% kapa: exam=1 round=33 year=2022 period=1 form=A subject=civil_law subject_ko=민법 seq=29 number=29
%% source: Q-Net 감정평가사 기출문제 (cst00309) artlSeq=5209997 file=raw/1008/5209997/2022년도 제33회 감정평가사 1차시험 1교시(민법,경제학,부동산학).pdf
%% conversion: kordoc 4.24.51 -> tools/gen1.py (auto); review=proofread-sonnet; flags=none
\begin{question}{29}
공유물 분할에 관한 설명으로 옳지 않은 것은? (다툼이 있으면 판례에 따름)
\begin{choices}
  \item 공유물분할청구권은 형성권에 해당한다.
  \item 공유관계가 존속하는 한 공유물분할청구권만이 독립하여 시효로 소멸될 수 없다.
  \item 부동산의 일부 공유지분 위에 저당권이 설정된 후 그 공유부동산이 현물분할된 경우, 저당권은 원칙적으로 저당권설정자에게 분할된 부분에 집중된다.
  \item 공유물분할청구의 소에서 법원은 원칙적으로 공유물분할을 청구하는 원고가 구하는 방법에 구애받지 않고 재량에 따라 합리적 방법으로 분할을 명할 수 있다.
  \item 공유자는 특별한 사정이 없는 한 언제든지 공유물의 분할을 청구할 수 있다.
\end{choices}
\end{question}
\answer{3}
